% !TeX root = dissertation.tex

\chapter{Zeroth-Order Mirror Descent}
\label{chap:zo-mirror-descent}

    Gradient descent iteravely takes a step in the gradient direction, gradually reducing function value. It successively builds better and better upper bounds on $f(x^*)$. However, it ignores the dual problem of finding lower bounds on $f(x^*)$. For convex functions, we have for all $u \in \R^d$,
    \begin{align*}
        f(u) \geq f(x) + \langle \grad f(x), u - x \rangle
    \end{align*}
    Thus, for some points $\set{x_t}_{t=0}^{T-1}$, we have
    \begin{align}
        \label{eq:bounded-by-regret}
        f(u) \geq \frac 1T \sum_{t=0}^{T-1} f(x_t) + \frac{1}{T} \sum_{t=0}^{T-1} \langle \grad f(x_t), u - x \rangle
    \end{align}
    Defining $\overline x = \frac{1}{T} \sum_{t=0}^{T-1} x_t$, convexity yields $f(\overline x) \leq \frac{1}{T} \sum_{t=0}^{T-1} f(x_t)$, and combining this with the above yields,
    \begin{align*}
        f(\overline x) - f(u) \leq \frac{1}{T} \sum_{t=0}^{T-1} \langle \grad f(x_t), x_t - u \rangle \eqqcolon R_T(u)
    \end{align*}
    $R_T(u)$, using terminology from online learing, is referred to as the regret. Consider  a reguarlized version 
    \begin{align*}
        \widetilde R_T(u) = \frac{1}{T} \left(-\frac{w(u)}{\alpha} + \sum_{t=0}^{T-1} \langle \grad f(x_t), x_t - u \rangle \right)
    \end{align*}
    where $\alpha > 0$ is a trade-off parameter and $w(u)$ is a strongly convex regularizer. By taking the updates $x_t = \mathrm{argmax}_u \widetilde R_T(u)$, we will tune $\alpha$ to show that Mirror descent converges in $T=O(\rho^2 / \ve^2)$ iterations where $\rho^2$ is the average value of $\mg{\grad f(x_t)}^2$ across iterations. This is because intuitively, the upper bound (\ref{eq:bounded-by-regret}) becomes better when $\frac{1}{T} \sum_{t=0}^{T-1} \mg{\grad f(x_t)}^2$ is small. 

    \begin{assumption}
        $f$ is $\rho$-Lipschitz continuous and convex over a closed convex set $Q \subset \R^d$.
    \end{assumption}

    \begin{definition}
        $w: Q \to \R^d$ is a DGF if it is 1-strongly convex w.r.t. $\mg{\cdot}$, in symbols, $w(y) \geq w(x) + \langle \grad w(x), y-x \rangle + \frac 12 \mg{y-x}^2$. The Bregman divergence is defined
        \begin{align*}
            V_x(y) = w(y) - \langle \grad w(y), y-x \rangle - w(x) 
        \end{align*}
        We have $V_x(x) = 0$ and $V_x(y) \geq \frac 12 \mg{x-y}^2 \geq 0$. 
    \end{definition}
    Classic examples are $w(x) = \frac 12 \mg{x}_2^2$ for the $\ell^2$ norm which gives $V_x(y) = \frac 12 \mg{y-x}_2^2$, and $w(y) = \sum_i y_i \log(y_i)$ for the $\ell^1$ norm with $Q = \Delta$, which gives KL divergence $V_x(y) = \sum_i y_i \log(y_i/x_i) \geq \frac 12 \mg{x-y}^2_1$. 

    \section{Classic Mirror Descent}
    \begin{definition}[Mirror Descent Step]
        The mirror descent step with step length $\alpha$ is 
        \begin{align*}
            x^+ = \Mirr_x(\alpha \grad f(x)) \quad \text{where} \quad \Mirr_x(\xi) = \underset{y \in Q}{\mathrm{argmin}} \set{V_x(y) + \langle \xi, y-x \rangle }.
        \end{align*}
    \end{definition}

    Mirror descents core lemma is this:
    \begin{lemma}[\cite{allenZhu2016linear}]
        \label{lem:mirror-descent}
        If $x_{t+1} = \Mirr_{x_t}(\alpha \grad f(x_k))$, then
        \begin{align*}
            \forall u \in Q,\quad \alpha (f(x_t) - f(u)) \leq \alpha \langle \grad f(x_t), x_t-u \rangle \leq \frac{\alpha^2}{2} \mg{\grad f(x_t)} + V_{x_t}(u) - V_{x_{t+1}}(u).
        \end{align*}
    \end{lemma}
    \begin{proof}
        Notice
        \begin{align*}
            \alpha \langle \grad f(x_t), x_t - u \rangle = \langle \alpha \grad f(x_t), x_t-x_{t+1} \rangle + \langle \alpha \grad f(x_t), x_{t+1} - u \rangle
        \end{align*}
        Since $Q$ is convex and $x_{t+1} = \Mirr_{x_t}(\alpha \grad f(x))$ is minimal, it satisfies the first order condition
        \begin{align*}
            \langle \grad V_{x_t}(x_{t+1}) + \alpha \grad f(x_t), u - x_{t+1} \rangle \geq 0 \quad \forall u \in Q
        \end{align*}
        Therefore the above is at most 
        \begin{align*}
            \langle \alpha \grad f(x_t), x_t - x_{t+1} \rangle + \langle - \grad V_{x_k}(x_{t+1}), x_{t+1} - u \rangle
        \end{align*}
        Using the three-point identity
        \begin{align*}
            \langle - \grad V_x(y) , y-u \rangle = V_x(u) - V_y(u) - V_x(y)
        \end{align*}
        we get the above equals
        \begin{align*}
            &\langle \alpha \grad f(x_t), x_t - x_{t+1} \rangle + V_{x_t}(u) - V_{x_{t+1}}(u) - V_{x_t}(x_{t+1}) \\
            \leq &\qty(\langle \alpha \grad f(x_t), x_t-x_{t+1} \rangle - \frac 12 \mg{x_t-x_{t+1}}^2) + (V_{x_t}(u) - V_{x_{t+1}}(u)) \\
            &\leq \frac{\alpha^2} 2 \mg{\grad f(x_t)}^2 + (V_{x_t}(u) - V_{x_{t+1}}(u))
        \end{align*}
        After recalling that $V_x(y) \geq \frac{1}{2}\mg{x-y}^2$, using Cauchy-Schwarz $\langle a, b \rangle \leq \mg{a} \cdot \mg{b}$ and Young's inequality $2ab \leq a^2 + b^2$ in that order.

        The first inequality follows because by convexity $f(u) \geq f(x_t) + \langle \grad f(x_t), u-x_t \rangle$.
    \end{proof}
    Using this, we can prove Mirror Descent converges.
    \begin{lemma}[\cite{allenZhu2016linear}]
        \label{lem:classic-tele}
        For $\bar x = \frac{1}{T} \sum_{t=0}^{T-1} f(x_t)$, if we choose $\alpha = \sqrt{2\Theta / \rho^2 T}$, we have
        \begin{align*}
            f(\bar x) - f(x^*) \leq \rho \sqrt{\frac{2\Theta }{T}}
        \end{align*}
    \end{lemma}
    \begin{proof}
        Telescoping, we get
        \begin{align*}
            \sum_{t=0}^{T-1} \alpha (f(x_t) - f(x^*)) \leq \sum_{t=0}^{T-1} \frac{\alpha^2}{2} \mg{\grad f(x_t)}^2 + V_{x_0}(x^*) - V_{x_T}(x^*)
        \end{align*}
        Using convexity, $f(\bar x) \leq 1/T \sum_{t=0}^{T-1} f(x_t)$, and recalling that $f$ is $\rho$-Lipschitz continuous, we have
        \begin{align*}
            f(\bar x) - f(x^*) \leq \frac{\alpha}{2} \rho^2 + \frac{1}{\alpha} \Theta,
        \end{align*}
        where $\Theta$ is any number satisfying $\Theta \geq V_{x_0}(x^*)$. Now by choosing $\alpha = \sqrt{2\Theta / \rho^2 T}$ to balance the two terms, we are done. 
    \end{proof}

    \section{Zero-Order Mirror Descent}
    In this section, we assume that the norm $\mg{\cdot}$ is the $\ell^2$ norm and $Q = \R^d$. We begin by introducing the zeroth-order gradient estimator.
    \begin{definition}[Zero-Order Estimator]
        The Zero-Order gradient estimator is defined as follows
        \begin{align*}
            \widetilde{\grad} f(x) = \frac{f(x+\sigma p) - f(x-\sigma p)}{2\sigma} \cdot p
        \end{align*}
        where $p \sim \mathcal N(0,I_d)$. 
    \end{definition}
    The main result of this chapter is the following.
    \begin{theorem}[informal]
        With iterates defined as $x_{t+1} = \Mirr_{x_t}(\alpha \widetilde \grad f(x_t))$, where the perturbation sampled at time $t$ is independent from the past, if $f$ is quadratic and
        \begin{align*}
            \alpha = \widetilde{O}\left(\sqrt{\frac{\Theta}{\rho^2 T d}}\right)
        \end{align*}
        if we let $\overline x = \frac{1}{T} \sum_{t=0}^{T-1} x_t$, with probability at least $1-\delta$, 
        \begin{align*}
            f(\overline x) - f^* \leq \widetilde{O}\left(\sqrt{\frac{d\Theta \rho^2}{T}}\right).
        \end{align*}
    \end{theorem}
    This matches the convergence rate of Mirror descent presented in~\cite{allenZhu2016linear}, up to the zero-order $d$ penalty and log factors.

    \section{Regret Bound}
    Now, notice that the only place in the proof of Lemma~\ref{lem:mirror-descent} where we used that the update is evaluated at the true gradient $\grad f(x_t)$ was to upper bound $f(x_t) - f(u)$. The gradient was not used in the proof, actually. So we get this lemma for free:
    \begin{lemma}
        If $x_{t+1} = \Mirr_{x_t}(\alpha \widetilde{\grad} f(x_t))$, then
        \begin{align*}
         \forall u \in Q \quad \alpha \langle \widetilde{\grad} f(x_t), x_t-u \rangle \leq \frac{\alpha^2}{2} \mg{\widetilde{\grad} f(x_t)}^2 + V_{x_t}(u) - V_{x_{t+1}}(u).
        \end{align*}
    \end{lemma}
    Notice now that for any $u \in Q$, 
    \begin{align*}
        \alpha (f(x_t) - f(u)) \leq \alpha \langle \grad f(x_t), x_t - u \rangle  = \alpha \langle \tilde \grad f(x_t), x_t - u \rangle + \alpha \langle \grad f(x_t) - \widetilde{\grad} f(x_t), x_t - u \rangle  
    \end{align*}
    
    For simplicity of exposition in this chapter only we will assume there is no quadratic error and $\widetilde{\grad} f(x_t) = p_tp_t^\top \grad f(x_t)$. Thus the above equals,
    \begin{align}
        &\alpha \langle \widetilde{\grad} f(x_t), x_t-u \rangle + \alpha \langle (I-p_tp_t^\top) \grad f(x_t), x_t-u \rangle \notag \\
        &\leq \frac{\alpha^2}{2} \mg{\widetilde{\grad} f(x_t)}^2 + V_{x_t}(u) - V_{x_{t+1}}(u) + \alpha \langle (I-p_tp_t^\top) \grad f(x_t), x_t-u \rangle \label{eq:telescope}
    \end{align}

    \section{Concentration Bounds}
    \begin{lemma}
        \label{lem:est-bound}
        With probability at least $1-\delta$, we have for every $1 \leq t \leq T$,
        \begin{align*}
            \mg{p_tp_t^\top \grad f(x_t)} \leq c \sqrt{d} \log(cT/\delta) \mg{\grad f(x_t)}
        \end{align*}
    \end{lemma}
    \begin{proof}
        By Lemma~\ref{lem:bounded-p-norm}, who's proof is deferred to Chapter~\ref{chap:fosp-eggroll}, the event
        \begin{align*}
            A_t = \set{\mg{p_t} \leq \sqrt{3d\log(C/\delta)}} 
        \end{align*}
        has probability at least $1-\delta$. Now, as $p_t^\top \grad f(x_t) \sim \mathcal N(0, \mg{\grad f(x_t)}^2)$ conditional on $\mathcal F_t$, we know that
        \begin{align*}
            \E[\exp(\frac{\mg{p_tp_t \grad f(x_t) \1_A}^2}{t^2})] \leq \E[\exp(\frac{3d \log(C/\delta) |p_t^\top \grad f(x_t)|^2}{t^2})]
        \end{align*}
        Using $\E[e^{\lambda Z^2}] = (1-2\lambda \sigma^2)^{-1/2}$ for $Z \sim \mathcal N(0, \sigma^2)$, we know that
        \begin{align*}
            \E[\exp(\frac{3d \log(C/\delta) |p_t^\top \grad f(x_t)|^2}{t^2}) \mid \mathcal F_t] = \left(1 - \frac{6d \log(C/\delta)\mg{\grad f(x_t)}^2}{t^2}\right)^{-1/2}
        \end{align*}
        The RHS is $\leq 2$ for $t = \sqrt{8d \log(C/\delta)}\mg{\grad f(x_t)}$. By taking expectations on both sides this bound also holds unconditionally. Then using Markov's inequality,
        \begin{align*}
            &\P(\mg{p_tp_t^\top \grad f(x_t) \1_A} \geq \sqrt{8d \log(C/\delta)}\mg{\grad f(x_t)} \lambda) 
            \\&= \P(\exp(\frac{\mg{p_tp_t^\top \grad f(x_t) \1_A}^2}{8d \log(C/\delta) \mg{\grad f(x_t)}^2}) \geq e^{\lambda^2}) \leq 2e^{-\lambda^2}
        \end{align*}
        Now choosing $\lambda = \sqrt{\log(C/\delta)}$, possibly changing $C$ to ensure $C \geq 2$, we have, with probability at least $1-\delta$ on $A$,
        \begin{align*}
            \mg{p_tp_t^\top \grad f(x_t)} \leq \log(C/\delta)\sqrt{8d}\mg{\grad f(x_t)} 
        \end{align*}
        Since $A$ has probability at least $1-\delta$, this bound holds with probability at least $1-2\delta$. Rescaling $\delta \gets \delta/2$ and changing the absolute constant yields the lemma.
    \end{proof}
    This lemma highlights the main technique for creating concentration bounds when the underlying distribution has fat tails. The first thing you do is truncate with high probability, and prove the resulting random variable has similar scale (whether that be sub-gaussian, sub-exponential, etc.) as what you want, up to log factors.


    The main idea behind the rest of this proof is that the mean-0 martingale error
    \begin{align*}
        \sum_{t=0}^{T-1} \alpha \langle (I-p_kp_k^\top) \grad f(x_k), x_k-u \rangle
    \end{align*}
    can hopefully be bounded by the other terms of the equation. In this case, it ends up being the ``energy'' term $V_{x_t}(u)$ and the Lipschitz constant $\rho$. 


    The main lemma which is different from the normal mirror descent bounds is the following.
    \begin{lemma}
        \label{lem:concentration_bounds}
    Let $\set{\mathcal F_t}$ be a filtration. Assume that $g_t, d_t \in \mathcal F_{t-1}$, and assume $p_t \sim \mathcal N(0,I)$ with $p_t \in \mathcal F_t$, $p_t \perp \mathcal F_{t-1}$. Define
    \begin{align*}
        E_\tau = \sum_{t=0}^{\tau-1} g_t^\top (I-p_tp_t^\top) d_t
    \end{align*}
    Then for any $\theta$, $0 < \delta < 1/e$, with probability at least $1-\delta$,
    \begin{align*}
        |E_\tau| \leq c\theta \sum_{t=0}^{\tau-1} \log(\frac{c\tau}{\delta}) \mg{g_t}^2 \mg{d_t}^2 + \frac1\theta \log(\frac {c\tau}{\delta}).
    \end{align*}
    \end{lemma}
    \begin{proof}
        For each $t$ let
        \begin{align*}
            A_t = \set{L_\delta \leq |p_t^\top d_t| \leq \sqrt{3 \log(c/\delta) \mg{d_t}^2}}, \quad \Pt(A_t) \geq 1-\delta
        \end{align*}
        Then by a lemma from before we know that $\Et[(I-p_tp_t^\top) d_t ; A_t] = 0$. In particular,
        \begin{align*}
            \Et[g_t^\top (I-p_tp_t^\top) d_t ; A_t] = 0.
        \end{align*} 
        Dropping the subscript $t$ for simplicity,
        \begin{align*}
            \E[\exp(\frac{(g^\top p p^\top d \1_A)^2}{t^2})] \leq \E[\exp(\frac{3 \log(C/\delta) \mg{d_t}^2(p^\top g)^2}{t^2})]
        \end{align*}
        Now, $p^\top g$ is $\mathcal N(0, \mg{g}^2)$ random variable. Therefore, $\E[\exp((p^\top g)^2 / 4\mg{g}^2)] \leq 2$. Now, taking $c \geq 1$ and $0 < \delta \leq 1/e$, we have
        \begin{align*}
            (g^\top (I-pp^\top) d)^2 \leq 2(g^\top d)^2 + 2 (g^\top p p^\top d)^2 
        \end{align*} 
        So, the random variable $Q_t = g_t^\top (I-p_tp_t^\top) d_t \1_{A_t}$ is $24\mg{g_t}^2 \mg{d_t}^2 \log(C/\delta)$ subgaussian and mean 0 given $\mathcal F_t$. Thus by the  (subgaussian) Hoeffding inequality with probability at least $1-\delta$ for any $\theta$,
        \begin{align*}
            \left |\sum_{t=0}^{\tau-1} Q_t  \right|\leq c \theta \sum_{t=0}^{\tau-1} \log(c/\delta) \mg{g_t}^2 \mg{d_t}^2 + \frac 1 \theta \log(\frac{c}{\delta})
        \end{align*}
        Now $\sum_{t=0}^{\tau-1} Q_t  = E_\tau$ on $\bigcap A_t$. Notice, since $\P(A_t) \geq 1-\delta$, $\P(\bigcap A_t) \geq 1-\tau \delta$. Therefore with probability at least $1-(\tau+1) \delta$, $E_\tau$ satisfies the same bound as above. Rescaling $\delta$ to $\delta/(\tau+1)$ and we have completed the proof.
    \end{proof}
    This above lemma is the key for making the argument work. Because $V_x(y) \geq \frac 12 \mg{x-y}^2$, and recalling that $f$ is $\rho$-Lipschitz continuous, we have
    \begin{align*}
        \left|\sum_{t=0}^{\tau-1} \alpha \langle (I-p_tp_t^\top) \grad f(x_t), x_t - u \rangle \right| \leq \alpha\sqrt{c \sum_{t=0}^{\tau-1} \log^2(c\tau/\delta) \rho^2 V_{x_t}(u)}
    \end{align*}
    after rescaling $c$ and by choosing the optimal $\theta$ according to AM-GM $(x+y)/2 \geq \sqrt{xy}$. So, it remains to bound $\sum V_{x_t}(u)$. We can do this because Lemma~\ref{lem:classic-tele} actually tells us a bound on $V_{x_\tau}(x^*)$. 

    \section{Energy Bound}
    So, by telescoping Equation~\eqref{eq:telescope}, we arrive at the following lemma.
    \begin{lemma}
        Define $\chi = \log(cT^2 / \delta)$. Then with probability at least $1-\delta$, for every $1 \leq \tau \leq T$,
        \begin{align*}
            V_{x_\tau}(x^*) \leq \frac{c \tau d \chi^{2} \alpha^2}{2} \rho^2 + V_{x_0}(x^*) + \alpha \rho \sqrt{c \sum_{t=0}^{\tau-1} \chi^{2} V_{x_t}(x^*)}
        \end{align*}
    \end{lemma}
\begin{proof}
    Telescoping Equation~\eqref{eq:telescope}, we get
    \begin{align*}
        \sum_{t=0}^{\tau-1} \alpha (f(x_t) - f(x^*)) &\leq \sum_{t=0}^{\tau-1} \frac{\alpha^2}{2} \mg{\tilde \grad f(x_t)}^2 + V_{x_0}(x^*) - V_{x_\tau}(x^*) 
        \\&\quad+ \sum_{t=0}^{\tau-1} \alpha \langle (I-p_tp_t^\top) \grad f(x_t), x_t - x^* \rangle 
    \end{align*}
    Therefore, by the previous discussion, by Lemma~\ref{lem:est-bound} and using that $f(x_t) \geq f(x^*)$ for any $x_t$,
    \begin{align*}
        V_{x_\tau}(x^*) \leq \frac{c\tau d \chi^{2} \alpha^2}{2} \rho^2 + V_{x_0}(x^*) + \alpha \rho \sqrt{c \sum_{t=0}^{\tau-1} \log^{2}(c\tau/\delta) V_{x_t}(x^*)}
    \end{align*}
    We can have this hold jointly for all $ 1 \leq \tau \leq T$ by rescaling $\delta \gets \delta/T$. Then by using $\tau \leq T$, we arrive at the statement of the lemma.
\end{proof}
Naturally, this above lemma should tell us how much $V_{x_t}(x^*)$ can grow. This next lemma tells us an upper bound on how much that is.
\begin{lemma}
    Let $\Theta$ be any upper bound on $V_{x_0}(x^*)$. Then for some absolute constant $c$,
    \begin{align*}
        \sup_{0 \leq t \leq T} V_{x_t}(x^*) \leq 2\Theta + cd\alpha^2 \rho^2  \chi^{2} T
    \end{align*}
\end{lemma}
\begin{proof}
    Define $M = \sup_{0 \leq t \leq T} V_{x_t}(x^*)$. Notice that the right hand side is increasing in $\tau$. Therefore applying sup to both sides, we get
    \begin{align*}
        M &\leq \frac{cT d \chi^{2} \alpha^2}{2} \rho^2 + \Theta + \alpha \rho \sqrt{c \sum_{t=0}^{T-1} \chi^{2} V_{x_t}(x^*)} \\
        &\leq \frac{cT d \chi^{2} \alpha^2}{2} \rho^2 + \Theta + \alpha \rho \sqrt{c T \chi^{2} M}
    \end{align*}
    Using Young's inequality, we know that
    \begin{align*}
        \alpha \rho \sqrt{cT\chi^{2} M} \leq \frac 12 M + c\alpha^2 \rho^2 \chi^{2} T 
    \end{align*}
    Therefore,
    \begin{align*}
        M &\leq cTd \chi^{2} \alpha^2 \rho^2 + 2\Theta + 2c\alpha^2 \rho^2 \chi^{2} T \\
        &\leq 2\Theta + cd\alpha^2 \rho^2 \chi^{2} T
    \end{align*}
    for some absolute constant $c$ as desired.
\end{proof}
A natural question is: Is this a good bound at all? Shouldn't $x_t$ be getting closer to the minmizer $x^*$? But we showed $V_{x_t}(x^*)$ grows linearly? 

To answer this, recall from before that $\alpha = \sqrt{2\Theta / \rho^2 T}$. Ignoring the polylog factor $\chi^{3/2}$, and dimension $d$, we have actually shown that $V_{x_t}(x^*)$ stays bounded. But still, shouldn't it be decreasing to 0? Actually, it should not, since $f$ could have many minimizers. The best bound you could hope to satisfy for every $f$ is that it is just bounded. 

For a similar choice of $\alpha$, this tells us the accumulated error 
\begin{align*}
    \sum_{t=0}^{T-1} \alpha \langle (I-p_tp_t^\top) \grad f(x_t), x_t-x^* \rangle
\end{align*}
is around the same size as $V_{x_0}(x^*)$. This observation is what let's us prove convergence of Zero-Order mirror descent.

\section{Putting it all Together}
\begin{theorem}[Zero-Order Mirror Descent Convergence]
    Let $\bar x = 1/T \sum_{t=0}^{T-1} x_t$, if we choose
    \begin{align*}
        \alpha = \sqrt{\frac{2\Theta}{\rho^2 c \chi^{2} T d}}
    \end{align*}
    then with probability at least $1-\delta$,
    \begin{align*}
        f(\bar x) - f(x^*) \leq \frac{1}{T\alpha} \cdot 9\Theta = \widetilde{O}\left(\rho \sqrt{\frac{\Theta d}{T}}\right)
    \end{align*}
\end{theorem}
\begin{proof}
    By telescoping and using $V_{x_T}(x^*) \geq 0$, we get
    \begin{align*}
        \sum_{t=0}^{T-1} \alpha (f(x_t) - f^*) \leq \frac{cTd \chi^{2} \alpha^2}{2} \rho^2 + \Theta + \alpha \sqrt{c \sum_{t=0}^{T-1} \chi^2 \rho^2 V_{x_t}(x^*)}
    \end{align*}
    By using the previous lemma, we know that
    \begin{align*}
        \alpha \sqrt{c \sum_{t=0}^{T-1} \chi^{2} \rho^2 V_{x_t}(x^*)} \leq \alpha \sqrt{cT\chi^{2} \rho^2(2\Theta + c \alpha^2 \rho^2 \chi^{2} T)}
    \end{align*}
    By choosing $\alpha = \sqrt{2\Theta / \rho^2 c \chi^{2} T d}\leq \sqrt{2\Theta / \rho^2 c \chi^{2} T}$ as $d \geq 1$, we know that
    \begin{align*}
        c \alpha^2 \rho^2 \chi^{2} T \leq 2\Theta
    \end{align*}
    and therefore,
    \begin{align*}
        \alpha \sqrt{cT\chi^{2} \rho^2(2\Theta + c \alpha^2 \rho^2 \chi^{2} T)} \leq 4\Theta
    \end{align*}
    Now again by our choice of $\alpha$,
    \begin{align*}
        \frac{cT d \chi^{2} \alpha^2}{2} \rho^2 \leq 4\Theta
    \end{align*}
    Therefore,
    \begin{align*}
        \alpha (f(\bar x) - f^*) \leq \frac{9\Theta}{T}
    \end{align*}
    We conclude,
    \begin{align*}
        f(\bar x) - f^* \leq \frac{9\Theta}{T \alpha} = \rho c \sqrt{\frac 92} \sqrt{\frac{d\Theta}{T}} \chi = \widetilde{O}\left(\rho \sqrt{\frac{d \Theta}{T}}\right).
    \end{align*}
\end{proof}

The extension of this anaslysis to the setting without quadratic error is left to future work. The main purpose of this chapter was not merely to show that mirror descent converges, but to determine whether zeroth-order mirror descent can be analyzed with the zeroth-order gradient estimator. Having answered this question affirmatively, the techniques developed here will be applied extensively in the next chapter.